\documentclass[11pt,a4paper]{article}

% ============================================================
% Packages (core only --- texlive-latex-base + texlive-base)
% ============================================================
\usepackage{amsmath, amssymb, amsthm}
\usepackage[a4paper, margin=2.5cm]{geometry}
\usepackage{array}

% ============================================================
% Theorem environments
% ============================================================
\theoremstyle{definition}
\newtheorem{definition}{Definition}[section]
\newtheorem{remark}{Remark}[section]
\newtheorem{proposition}{Proposition}[section]
\newtheorem{assumption}{Assumption}[section]

\theoremstyle{plain}
\newtheorem{lemma}{Lemma}[section]

% ============================================================
% Notation macros
% ============================================================
\newcommand{\R}{\mathbb{R}}
\newcommand{\E}{\mathbb{E}}
\newcommand{\N}{\mathbb{N}}
\newcommand{\Var}{\operatorname{Var}}
\newcommand{\Cov}{\operatorname{Cov}}
\newcommand{\norm}[1]{\left\lVert #1 \right\rVert}
\newcommand{\inner}[2]{\langle #1,\, #2 \rangle}
\newcommand{\corr}{\operatorname{corr}}
\newcommand{\sgn}{\operatorname{sign}}
\DeclareMathOperator{\IC}{IC}
\DeclareMathOperator{\IR}{IR}
\DeclareMathOperator{\MI}{I}
\DeclareMathOperator*{\argmax}{arg\,max}
\newcommand{\code}[1]{\texttt{\detokenize{#1}}}

% ============================================================
\begin{document}

\title{\textbf{From Narrative to Alpha}\\[4pt]
\large Exogenous Perception Pressure as an Orthogonal Signal in an\\
Autonomous Microstructure Alpha Book: Timescale Separation,\\
Effective Breadth, and the Adverse-Selection Barrier}

\author{Yigit Onat\\
\textit{Independent Researcher}\\
\texttt{yionat@gmail.com}}
\date{July 2026}
\maketitle

% ============================================================
\begin{abstract}
A book of alpha signals built entirely from order-book microstructure is powerful but structurally
narrow: its signals share one information source (order flow), occupy one band of timescales
(sub-second to minutes), and are jointly capped by an adverse-selection barrier that collapses a
large unconditional information coefficient once one conditions on the states in which a passive fill
actually occurs. We ask what a \emph{single exogenous signal from a different information domain}
contributes to such a book, and we answer at the level of methodology rather than measured
performance. Taking \emph{Perception Pressure} $P(K,t)$---a signed, credibility-weighted media-narrative
signal produced by the Prism system---as the exogenous input, we specify (i) a source-agnostic
integration architecture in which a narrative series is aligned to market bars by a point-in-time
as-of join and scored by the same information-coefficient, mutual-information, and five-gate
replication machinery as any endogenous feature, and (ii) three theoretical reasons the addition can
improve the book that do not require the exogenous signal to be individually strong. First,
\emph{timescale separation}: narrative momentum lives in a spectral band essentially disjoint from the
$0.005$--$0.1$\,Hz band that carries order-book-imbalance predictability, so the two are close to
orthogonal by construction. Second, \emph{effective breadth}: under the fundamental law of active
management the information ratio scales with the effective number of independent bets
$N_{\mathrm{eff}}=N/[1+(N-1)\bar\rho]$, and appending a low-correlation signal raises
$N_{\mathrm{eff}}$ more than appending another correlated microstructure signal of equal strength.
Third, \emph{partial evasion of the adverse-selection barrier}: because a slow exogenous signal is
harvested by patient execution over its own horizon rather than by reactive spread-crossing, its
coupling to the fillable-state indicator---the mechanism that destroys microstructure edge---is
weaker, so its conditional information coefficient need not collapse in the same way. Each claim is
stated as a hypothesis with an explicit, leakage-controlled, false-discovery-rate-corrected test; we
report no performance figures. The contribution is a theory and a protocol for making an autonomous
microstructure research platform \emph{more complete}---a decorrelated, multi-domain alpha book that
reports its own honest effective breadth---by admitting narrative as a first-class, gated signal
source.
\end{abstract}

\smallskip
\noindent\textbf{Keywords:} exogenous signal, alpha combination, signal orthogonality, effective
breadth, fundamental law of active management, adverse selection, mutual information, information
coefficient, narrative, market microstructure

\smallskip
\noindent\textbf{JEL Classification:} C55, C58, G14, G17

\tableofcontents
\bigskip

% ============================================================
\section{Introduction}
\label{sec:intro}
% ============================================================

\subsection{The completeness problem of a single-domain alpha book}

An autonomous microstructure research platform can, in principle, manufacture an unbounded number of
features from the order book: imbalances, depth slopes, queue dynamics, trade-sign runs, spectral and
information-theoretic transforms. In practice these features are far from independent. They are
computed from the same underlying object---the limit order book and the trade tape---so they share
information, co-move, and, when combined, deliver far less than the naive sum of their individual
information coefficients. Three structural facts bound such a book:

\begin{enumerate}
  \item \textbf{One information source.} Every signal is a functional of order flow. Whatever order
        flow cannot see---an exogenous shock forming in public discourse before it reaches the
        tape---no such signal can anticipate.
  \item \textbf{One band of timescales.} Order-book-imbalance predictability is empirically
        concentrated at high frequency and decays to noise within minutes; a companion study on the
        same venue localizes essentially all of its predictive power to the $0.005$--$0.1$\,Hz band
        (periods of $10$--$200$\,s) with an Ornstein--Uhlenbeck mean-reversion half-life of a few
        seconds. A book living only in this band is exposed to a single regime of decay.
  \item \textbf{A shared economic ceiling.} On a decentralized perpetual-futures venue, a strong
        \emph{unconditional} predictor can be economically inert: conditioning the information
        coefficient on the states in which a directionally-correct passive fill would actually occur
        collapses it from large to near zero. This adverse-selection barrier is a property of order
        flow itself and therefore afflicts order-flow signals \emph{jointly}.
\end{enumerate}

A book that is deep in one domain but empty in all others is thus fragile in a precise sense: its
signals fail together, decay together, and are filled against together. Robustness---and additional,
genuinely new information ratio---requires a signal that is \emph{not} a function of order flow.

\subsection{Narrative as the exogenous domain}

The natural exogenous domain for a liquid, sentiment-driven asset class is public narrative. We take
as given the \emph{Perception Pressure} signal $P(K,t)$ and its momentum $\dot P(K,t)$ defined in the
companion instrument paper (Onat, 2026a): a per-concept, signed, credibility-weighted media-narrative
series, sampled on a fixed cadence, in which attention and directional sentiment are multiplicatively
coupled and no single source can manufacture the measure. Our concern here is not how $P$ is
computed---that is the instrument paper's subject---but what it is \emph{worth to a trading book} and
how to admit it without contaminating an otherwise disciplined research process.

\subsection{Contributions}

\begin{enumerate}
  \item \textbf{A source-agnostic integration architecture} (Section~\ref{sec:arch}) in which an
        exogenous narrative series is aligned to market bars by a strictly backward-looking as-of join
        and then scored---by information coefficient across horizons, by mutual information, and
        through a five-gate replication protocol---by exactly the machinery that scores an endogenous
        feature. The narrative signal is not privileged; it earns its place by the same gates.
  \item \textbf{A timescale-separation argument} (Section~\ref{sec:timescale}) that narrative momentum
        and order-book-imbalance predictability occupy near-disjoint spectral bands, making the
        exogenous signal approximately orthogonal to the existing book by construction rather than by
        luck.
  \item \textbf{An effective-breadth argument} (Section~\ref{sec:breadth}) that formalizes, via the
        fundamental law of active management and the effective number of independent bets, why a
        low-correlation signal contributes more marginal information ratio than another correlated
        microstructure signal of equal individual strength---and why the binding quantity is
        \emph{stress-regime} correlation.
  \item \textbf{An adverse-selection argument} (Section~\ref{sec:advsel}) that a slow exogenous signal,
        harvested by patient execution over its own horizon, has weaker coupling to the fillable-state
        indicator than a fast order-flow signal, and therefore need not suffer the same conditional-IC
        collapse.
  \item \textbf{A falsifiable evaluation protocol} (Section~\ref{sec:eval}) stating, for each claim, a
        leakage-controlled and multiple-testing-corrected test with explicit kill criteria. We report
        no empirical figures; this is a methods paper.
\end{enumerate}

We are explicit throughout that the value of an exogenous signal is an empirical question. What this
paper supplies is the theory that says \emph{where} to look for that value, and the protocol that
keeps the search honest.

% ============================================================
\section{Background and Related Work}
\label{sec:related}
% ============================================================

\paragraph{Signal combination and breadth.}
Grinold's fundamental law of active management (Grinold, 1989; Grinold and Kahn, 2000) relates the
information ratio of a strategy to the skill of its signals and the number of independent bets,
$\IR \approx \IC\cdot\sqrt{\mathrm{BR}}$. When signals are correlated the naive breadth $\mathrm{BR}=N$
overstates the achievable ratio; the operative quantity is an \emph{effective} breadth that discounts
for redundancy. We use the effective-number-of-bets form
$N_{\mathrm{eff}}=N/[1+(N-1)\bar\rho]$ developed for this platform's process framework (Onat, 2026c),
which makes the marginal value of a new signal an explicit function of its correlation to the
incumbent book.

\paragraph{Adverse selection and the limits of microstructure alpha.}
That a statistically strong microstructure predictor can be economically inert is documented for this
venue in the companion microstructure-alpha study (Onat, 2026b), where the imbalance information
coefficient of $\approx 0.45$ at second horizons collapses to $\approx 0.03$ once restricted to
fillable states---an empirical instance of the classical adverse-selection problem of informed versus
uninformed order flow (Glosten and Milgrom, 1985; Kyle, 1985). The barrier motivates the search for
signals whose exploitation does not require winning a race for the same passive fills.

\paragraph{News, attention, and returns.}
That media tone and attention carry return-relevant information is long established (Tetlock, 2007;
Da, Engelberg, and Gao, 2011); the instrument paper (Onat, 2026a) situates Perception Pressure within
this literature. The present paper is downstream of all of it: we treat the existence of \emph{some}
narrative-return relationship as a hypothesis to be tested under controls, and study what its
inclusion does to a microstructure book's \emph{portfolio} of signals.

\paragraph{Information-theoretic evaluation.}
We adopt mutual information as the currency of predictive content, justified through Fano's inequality
as a bound on achievable prediction error (Cover and Thomas, 2006), estimated by the
$k$-nearest-neighbour method (Kraskov, St\"ogbauer, and Grassberger, 2004), and controlled for
multiple testing by the Benjamini--Hochberg procedure (Benjamini and Hochberg, 1995). These are
imported from the consuming platform, not invented here.

% ============================================================
\section{The Exogenous Signal}
\label{sec:signal}
% ============================================================

We recall only what is needed downstream. For a tracked concept $K$ (an entity mapped to a tradable
symbol, e.g.\ ``Bitcoin'' $\to$ \code{BTC}), the instrument produces on a fixed cadence a salience
$A(K,t)\ge 0$, a valence $V(K,t)\in[-1,1]$, a signed perception $P(K,t)=A(K,t)\,V(K,t)$, and its
momentum $\dot P(K,t)=P(K,t)-P(K,t-1)$.

\begin{assumption}[Entity map]
\label{ass:map}
A fixed, time-invariant, publicly reproducible map $\phi$ sends tradable symbols to sets of tracked
concepts, $\phi(\code{BTC})=\{K_1,\dots\}$. When several concepts map to one symbol their signals are
aggregated by a salience-weighted mean, itself computed only from information available at $t$.
\end{assumption}

\begin{remark}[Momentum is the leading candidate]
\label{rem:momentum}
As argued in the instrument paper, a level $P$ that merely co-moves with price is contemporaneous and
of limited predictive use, whereas a narrative that is \emph{building}, $\dot P>0$, before it is fully
priced is the economically interesting case. We therefore treat $\dot P$ (and higher-order changes) as
the primary candidate and $P$, $A$, $V$ as secondary. Nothing in the architecture depends on this
choice; it is a prior about where predictive content is most likely, to be confirmed or rejected by
the gates.
\end{remark}

% ============================================================
\section{Integration Architecture}
\label{sec:arch}
% ============================================================

The design principle is that the exogenous signal must be \emph{indistinguishable}, to the scoring
machinery, from an endogenous feature---so that it is held to the same standard and enjoys no special
pleading.

\subsection{The feature-source contract}

A narrative series enters through a minimal contract: a source exposes
\code{load(symbol, start, end)} returning a data frame indexed by a UTC timestamp with one column per
narrative component. The platform's process runner merges these columns onto market bars and, from
that point on, a narrative column is scored by \code{ic_horizon} and \code{mi_ksg} exactly as an
order-book column is. Processes operate on columns and never learn a column's origin. This is the
formal sense in which narrative is a first-class signal: it is admitted by the type system, not by a
bespoke code path.

\subsection{Point-in-time alignment}

\begin{definition}[Backward as-of join]
\label{def:asof}
Let market bars close at times $\{t_i\}$ and let narrative snapshots be stamped at $\{\tau_j\}$, each
$\tau_j$ being the instant the snapshot became \emph{queryable}. The aligned value on bar $i$ is
$P^{\star}(t_i)=P(\tau_{j^\star})$ with $j^\star=\max\{j:\tau_j\le t_i-\delta\}$, where $\delta\ge 0$ is
a mandatory latency guard.
\end{definition}

The join is strictly backward-looking and carries an explicit latency guard $\delta$ absorbing
ingestion, analysis, and publication delay. Two failure modes are thereby excluded by construction:
using a snapshot timestamped before it was computable (\emph{look-ahead}), and using the revised value
of a snapshot later corrected (\emph{restatement}); the source must expose the value as first known,
not as later amended.

\subsection{Reusing the five-gate replication protocol}

Once aligned, a narrative column traverses the same protocol as any candidate (Onat, 2026b):
(1) a \emph{discovery} gate testing information coefficient against forward returns across horizons
under Benjamini--Hochberg control; (2) a \emph{cost} gate re-testing net of realistic fees, funding,
and slippage; (3) a \emph{temporal} gate requiring persistence across disjoint time folds;
(4) a \emph{cross-symbol} gate requiring the effect to replicate across entities that map to distinct
concepts; and (5) a \emph{correlation} gate testing incremental content against the incumbent book.
A narrative signal that fails any gate is discarded exactly as an endogenous one would be. Crucially,
the horizons at which narrative is tested are longer than those for microstructure---hours to days,
not seconds---because that is where Section~\ref{sec:timescale} predicts its content lies; the gate
thresholds and FDR ledger are shared, the horizon grid is not.

% ============================================================
\section{Why It Helps I: Timescale Separation}
\label{sec:timescale}
% ============================================================

\begin{proposition}[Near-orthogonality by band disjointness]
\label{prop:band}
Let $x(t)$ be an order-book-imbalance signal whose predictive cross-spectrum with returns is supported
on a high-frequency band $B_x$ (empirically $\approx[0.005,0.1]$\,Hz on this venue), and let $\dot P(t)$
be narrative momentum sampled at cadence $\Delta$ with $1/\Delta \ll \min B_x$, so that its spectral
content is supported on a low-frequency band $B_P$ with $B_P\cap B_x=\varnothing$. If both are
approximately band-limited to their supports, their contemporaneous correlation obeys
\[
  |\corr(x,\dot P)| \;\le\; \frac{\int_{B_x\cap B_P} \sqrt{S_x(f)\,S_{\dot P}(f)}\,df}
       {\sqrt{\int S_x \, \int S_{\dot P}}} \;\approx\; 0,
\]
by the Cauchy--Schwarz bound on cross-spectral density over disjoint supports.
\end{proposition}

The content of Proposition~\ref{prop:band} is not that narrative and microstructure are provably
uncorrelated---real signals are only approximately band-limited, and spectral leakage leaves a small
residual---but that their correlation is \emph{structurally suppressed}: they carry information at
different frequencies, so a book combining them is combining near-independent bets rather than
repackaging the same one. This is precisely the property that the effective-breadth argument rewards.

\begin{remark}[Why the venue's spectral result matters]
\label{rem:spannung}
The empirical localization of imbalance predictability to a narrow high-frequency band (companion
spectral study) is what makes the disjointness concrete: it is not merely that narrative is ``slower''
in the abstract, but that the incumbent book's information demonstrably vanishes above periods of a
few hundred seconds, exactly where a narrative signal sampled every few minutes to hours begins.
\end{remark}

% ============================================================
\section{Why It Helps II: Effective Breadth}
\label{sec:breadth}
% ============================================================

We formalize the marginal value of a low-correlation signal. Consider a book of $N$ microstructure
signals, each standardized, with common individual information coefficient $\IC$ and common pairwise
correlation $\rho\in[0,1)$. Its effective breadth is
\begin{equation}
  N_{\mathrm{eff}}(N,\rho) = \frac{N}{1+(N-1)\rho},
  \label{eq:neff}
\end{equation}
and by the fundamental law $\IR\approx \IC\sqrt{N_{\mathrm{eff}}}$.

\begin{proposition}[A low-correlation signal buys more breadth]
\label{prop:marginal}
Append one further signal of equal individual strength $\IC$ that has correlation $\rho_x$ with each
incumbent. Write the augmented effective breadth as $N_{\mathrm{eff}}'(\rho_x)$ computed from the
$(N{+}1)\times(N{+}1)$ correlation matrix with incumbent block correlation $\rho$ and new-row
correlation $\rho_x$. Then $N_{\mathrm{eff}}'$ is strictly decreasing in $\rho_x$; consequently an
exogenous signal with $\rho_x\approx 0$ (Proposition~\ref{prop:band}) delivers strictly greater
marginal effective breadth than an additional microstructure signal with $\rho_x=\rho>0$.
\end{proposition}

\begin{proof}[Sketch]
The mean off-diagonal correlation of the augmented book is
$\bar\rho' = \big[N(N-1)\rho + 2N\rho_x\big]/\big[(N+1)N\big]$, which is increasing in $\rho_x$; the
average-correlation form of the effective breadth,
$N_{\mathrm{eff}}'=(N+1)/[1+N\bar\rho']$, is decreasing in $\bar\rho'$ and hence in $\rho_x$. The
final claim is the comparison at $\rho_x\approx 0$ versus $\rho_x=\rho$.
\end{proof}

Proposition~\ref{prop:marginal} is the quantitative statement of ``more complete.'' Beyond some depth,
the $N{+}1$-th microstructure signal is nearly redundant---it raises $N$ but also $\bar\rho$, and
$N_{\mathrm{eff}}$ barely moves---whereas the first signal from a new, near-orthogonal domain moves
$N_{\mathrm{eff}}$ by close to a full unit even at modest individual $\IC$. A book's information ratio
is grown far more cheaply by \emph{widening its domains} than by \emph{deepening one domain}.

\begin{remark}[Stress-regime correlation is the binding quantity]
\label{rem:stress}
The correlations in \eqref{eq:neff} must be measured where diversification is actually needed: in
stress. Signals that are decorrelated in calm markets frequently co-move when volatility spikes,
collapsing realized breadth exactly when it is demanded. A narrative signal is a candidate to break
this pattern precisely because its driver---public discourse---is not the order flow that convulses in
a microstructure stress event; but this is a hypothesis, and the evaluation therefore estimates
$\rho_x$ \emph{conditionally on volatility regime}, not marginally.
\end{remark}

% ============================================================
\section{Why It Helps III: The Adverse-Selection Barrier}
\label{sec:advsel}
% ============================================================

The sharpest limit on microstructure alpha is not weak prediction but adverse selection: the signal
predicts the future, yet a directionally-correct passive order is filled preferentially in the states
where the prediction is wrong. Formally, for signal $g$ and forward return $r$, let $F$ be the event
that a passive order placed on $g$'s side is filled. The exploitable skill is the \emph{conditional}
information coefficient $\IC(g,r\mid F)$, and the barrier is the gap
$\IC(g,r)-\IC(g,r\mid F)$, which the companion study finds to be almost the whole of $\IC$ for
order-book imbalance.

\begin{proposition}[Slow signals couple weakly to the fill state]
\label{prop:advsel}
Suppose signal $g$ is exploited not by a reactive passive quote at the touch but by a schedule that
works into the position over a horizon $H_g$ comparable to $g$'s own decay timescale. Then the
relevant conditioning event is not the instantaneous touch-fill $F$ but the average execution outcome
over $[t,t+H_g]$, whose dependence on the fast fill-selection mechanism is diluted by a factor growing
with $H_g$ relative to the spread-crossing timescale $\tau_s$. Hence the conditional-IC degradation of
a signal with $H_g\gg\tau_s$ is bounded above by that of a signal with $H_g\approx\tau_s$, and can be
strictly smaller.
\end{proposition}

The intuition is executional. Microstructure edge at second horizons must be captured by winning
passive fills in a race, and the winner of that race is systematically the adversely-selected side. A
narrative signal with an hours-to-days horizon is captured differently: one has time to work orders
patiently, cross spreads selectively, and split execution, so the pathological correlation between
``getting filled'' and ``being wrong'' is materially weaker. Proposition~\ref{prop:advsel} does not
claim the barrier vanishes---exogenous signals face their own frictions, including that a widely-read
narrative may already be priced---but it identifies a concrete channel by which a slow signal can
convert a larger fraction of its statistical content into economic content than a fast one.

% ============================================================
\section{Evaluation Protocol}
\label{sec:eval}
% ============================================================

Each claim above is falsifiable; we specify the test and its kill criterion, and report no numbers.

\begin{enumerate}
  \item \textbf{Orthogonality (tests Prop.~\ref{prop:band}).} Estimate the volatility-regime-conditional
        correlation $\rho_x$ between aligned narrative momentum and the incumbent imbalance signal, and
        the cross-spectral coherence across frequency. \emph{Kill:} if $\rho_x$ is not small in stress
        regimes, the effective-breadth argument does not apply and the signal offers no diversification.
  \item \textbf{Incremental predictive content (tests the reason to include at all).} Estimate the
        conditional mutual information $\MI(r_{t\to t+h};\,\dot P_t \mid \text{price}_t,\text{volume}_t)$
        at narrative-scale horizons $h$, under the $k$NN estimator with a block permutation null and
        Benjamini--Hochberg control across concepts and horizons. \emph{Kill:} if conditional MI is not
        significantly positive after FDR, narrative adds nothing beyond price and volume.
  \item \textbf{Breadth realization (tests Prop.~\ref{prop:marginal}).} Compute the book's
        $N_{\mathrm{eff}}$ and out-of-sample information ratio with and without the narrative signal,
        walk-forward. \emph{Kill:} if $N_{\mathrm{eff}}$ and OOS IR do not rise on inclusion, the
        marginal-breadth claim fails in practice.
  \item \textbf{Execution realizability (tests Prop.~\ref{prop:advsel}).} Under a schedule-based
        execution model at horizon $H_g$, estimate the conditional information coefficient
        $\IC(\dot P,r\mid F)$ and compare its degradation to that of the imbalance signal. \emph{Kill:}
        if the conditional IC collapses as severely as for microstructure, the adverse-selection
        advantage is illusory.
  \item \textbf{Leakage and stability audit.} Re-run studies 1--4 with the latency guard $\delta$
        inflated and with the entity map $\phi$ frozen strictly before the test window. \emph{Kill:} if
        results depend on a short $\delta$ or on post-hoc concept selection, they are artifacts.
\end{enumerate}

Study~2 is the necessary condition---without incremental information there is nothing to integrate---and
studies~3 and~4 are what distinguish a signal that merely predicts from one that improves a
\emph{book}. All five are framed as leakage-controlled and multiple-testing-corrected because a signal
tested against many concepts and horizons is otherwise a textbook data-snooping trap.

% ============================================================
\section{What ``More Complete'' Means}
\label{sec:complete}
% ============================================================

The object this paper argues for is not a narrative strategy but a \emph{multi-domain, decorrelated
alpha book that reports its own honest effective breadth}. In that book the microstructure signals
remain the high-frequency core; the narrative signal is a low-frequency, near-orthogonal overlay
admitted through the same gates and deployed as one more process, not as a parallel system. Three
consequences follow. The book's information ratio can grow through width rather than through ever more
correlated depth (Section~\ref{sec:breadth}). Its failure modes diversify: a microstructure regime
break need not coincide with a narrative regime break (Remark~\ref{rem:stress}). And a portion of its
statistical content may survive the executional filter that caps its order-flow signals
(Section~\ref{sec:advsel}). Sophistication here is not additional complexity for its own sake but the
disciplined admission of a second information domain under an unchanged standard of proof---the
platform becomes more complete by knowing, and gating, what it does not otherwise see.

% ============================================================
\section{Limitations}
\label{sec:limitations}
% ============================================================

We disclose the principal threats. The narrative signal inherits every measurement limitation of its
instrument (Onat, 2026a): language-model sentiment error, lexical rather than semantic clustering, and
frequently-inert engagement terms; these bound the fidelity of $\dot P$ and hence of any measured
effect. Narrative events are \emph{sparse} relative to microstructure ticks, so statistical power at
narrative horizons is limited and the multiple-testing correction is correspondingly punishing---an
honest evaluation may simply fail to reject the null for want of events. The narrative--return
relationship is plausibly \emph{non-stationary}: a signal that leads price while under-attended can
decay to contemporaneous as it becomes widely read, so any positive finding requires ongoing decay
monitoring rather than a one-time validation. The entity map $\phi$ is a researcher choice and a
degree of freedom that must be frozen before testing. Finally, the theoretical propositions are
deliberately modest: Proposition~\ref{prop:band} bounds correlation only for approximately
band-limited signals, Proposition~\ref{prop:marginal} assumes equal individual skill for its clean
comparison, and Proposition~\ref{prop:advsel} identifies a channel without quantifying it. None of
these is a claim of profit; each is a signpost for the evaluation.

% ============================================================
\section{Conclusion}
\label{sec:conclusion}
% ============================================================

We have set out how, and why, to admit an exogenous narrative signal into an autonomous microstructure
alpha book. The how is an integration architecture that renders narrative indistinguishable from an
endogenous feature to the scoring machinery, aligned point-in-time and gated by an unchanged five-gate
protocol. The why is three structural arguments---timescale separation making the signal near-orthogonal
by construction, effective breadth rewarding that orthogonality more than additional correlated depth,
and a weaker coupling to the fill state partially evading the adverse-selection barrier that caps
order-flow alpha. Each is a hypothesis with a stated, leakage-controlled, FDR-corrected test and an
explicit kill criterion; none is a performance claim. The contribution is a principled route to a more
complete quantitative research platform: not a bigger pile of order-book features, but a second
information domain, admitted under the same standard of proof, whose value---if it exists---shows up
exactly as increased effective breadth and surviving conditional skill. Turning the protocol on data
is the next step.

% ============================================================
\section*{References}
\addcontentsline{toc}{section}{References}
% ============================================================

\begin{enumerate}
  \item Benjamini, Y.\ \& Hochberg, Y.\ (1995). Controlling the false discovery rate: a practical and
        powerful approach to multiple testing. \textit{Journal of the Royal Statistical Society,
        Series B}, 57(1), 289--300.

  \item Cover, T.M.\ \& Thomas, J.A.\ (2006). \textit{Elements of Information Theory} (2nd ed.).
        Wiley-Interscience.

  \item Da, Z., Engelberg, J.\ \& Gao, P.\ (2011). In search of attention.
        \textit{Journal of Finance}, 66(5), 1461--1499.

  \item Glosten, L.R.\ \& Milgrom, P.R.\ (1985). Bid, ask and transaction prices in a specialist
        market with heterogeneously informed traders. \textit{Journal of Financial Economics}, 14(1),
        71--100.

  \item Grinold, R.C.\ (1989). The fundamental law of active management.
        \textit{Journal of Portfolio Management}, 15(3), 30--37.

  \item Grinold, R.C.\ \& Kahn, R.N.\ (2000). \textit{Active Portfolio Management} (2nd ed.).
        McGraw-Hill.

  \item Kraskov, A., St\"ogbauer, H.\ \& Grassberger, P.\ (2004). Estimating mutual information.
        \textit{Physical Review E}, 69(6), 066138.

  \item Kyle, A.S.\ (1985). Continuous auctions and insider trading. \textit{Econometrica}, 53(6),
        1315--1335.

  \item Onat, Y.\ (2026a). Perception Pressure and Resonance: decomposable, trust-weighted measures of
        media narrative from a multi-perspective news-aggregation system. Working paper (Prism).

  \item Onat, Y.\ (2026b). Autonomous microstructure alpha discovery under false-discovery control:
        order-book imbalance, conditional information coefficients, and the adverse-selection barrier
        in Hyperliquid perpetual futures. Working paper.

  \item Onat, Y.\ (2026c). The Process: a first-class analytical unit for orthogonal information
        extraction from market microstructure. Working paper.

  \item Tetlock, P.C.\ (2007). Giving content to investor sentiment: the role of media in the stock
        market. \textit{Journal of Finance}, 62(3), 1139--1168.
\end{enumerate}

\end{document}
